\documentclass[10pt,a4paper]{article}
\usepackage{amsmath,amssymb,amsthm,graphicx}
\title{Symmetry of time-bounded information from meta-complexity algorithms and average-case easiness}
\author{Shuichi Hirahara}
\date{doi:10.4230/LIPIcs.CCC.2022.26; CC BY4.0}
\begin{document}
\maketitle
\noindent\textit{Editorial note.} Original author: Shuichi Hirahara. Official CCC2022 article26 (2022-07-11), DOI 10.4230/LIPIcs.CCC.2022.26. Creative Commons Attribution4.0.\par\medskip
\noindent\textit{Editorial note.} License: https://creativecommons.org/licenses/by/4.0/\par\medskip
\noindent\textit{Editorial note.} Faithful editable TeX transcription of formal Theorem4.1 and author main Theorem1.2, complete Section4 proof pp26:15-26:16, exact necessary definitions and cited prior tools from Sections2.1/3, and complete published bibliography. Original author mathematical expressions and numbering retained. Standalone typography and long-equation line breaks adapted; no new proofs or mathematical corrections.\par\medskip
\noindent\textit{Allowed standard prerequisites (editorial).}\begin{itemize}
\item Approximation of time-bounded Kolmogorov complexity from errorless NP easiness. Exact author Lemma2.1 citing prior [31,33] retained.
\item Deterministic direct-product reconstruction under an E circuit lower bound. Exact prior author Lemma2.2 citing [35], generator definition, and Fact3.4 retained.
\end{itemize}\medskip
\paragraph{3. Preliminaries: Notation.}
$[n]$ denotes $\{1,\ldots,n\}$. For a function $p:\mathbb N\to\mathbb N$ and $i\in\mathbb N$, the function $p^{(i)}:\mathbb N\to\mathbb N$ is recursively defined as follows: $p^{(0)}(n):=n$ and $p^{(i+1)}(n):=p^{(i)}(p(n))$ for every $n\in\mathbb N$.
\paragraph{Kolmogorov Complexity.}
For a string $x\in\{0,1\}^*$ and an index $i\in\mathbb N$, let $x_i$ denote the $i$-th bit of $x$ if $i\leq|x|$ and $\bot$ if $i>|x|$. We fix an efficient universal Turing machine $U$. Time-bounded Kolmogorov complexity is formally defined as follows.
\paragraph{Definition 3.1 (Time-bounded Kolmogorov complexity).}
For strings $x,y\in\{0,1\}^*$, a time bound $t\in\mathbb N\cup\{\infty\}$, and an oracle $A$, the $A$-oracle $t$-time-bounded Kolmogorov complexity of $x$ given $y$ is defined as
\[K^{t,A}(x\mid y):=\min\{|d|:U^{A,y,d}\text{ outputs }x_i\text{ on input }i\text{ in time }t\text{ for every }i\in[1,\ldots,|x|+1]\}.\]
Here, $U^{A,y,d}$ indicates that the universal Turing machine is given oracle access to $A$ and each bit of $y$ and $d\in\{0,1\}^*$. We omit the superscript $A$ if $A=\varnothing$, the superscript $t$ if $t=\infty$, and ``$\mid y$'' if $y$ is the empty string.
A simple counting argument implies the following basic fact of Kolmogorov-randomness.
\paragraph{Fact 3.2.}
For any integer $s\geq1$ and any string $y\in\{0,1\}^*$, the number of strings $x\in\{0,1\}^*$ such that $K(x\mid y)<s$ is less than $2^s$.
\begin{proof}
The number of programs of length less than $s$ is at most $\sum_{i=0}^{s-1}2^i<2^s$.
\end{proof}
\paragraph{Definition 3.3 ([50,33]).}
For a polynomial $\tau:\mathbb N\to\mathbb N$, let
\[\Pi_{\mathrm{Yes}}:=\{(x,1^t,1^s):K^t(x)\leq s\},\]
\[\Pi_{\mathrm{No}}:=\{(x,1^t,1^s):K^{\tau(|x|,t)}(x)>s+\log\tau(|x|,t)\}.\]
We define $\mathsf{Gap}_{\tau}\mathsf{MINKT}$ to be the promise problem $(\Pi_{\mathrm{Yes}},\Pi_{\mathrm{No}})$. We say that $\mathsf{GapMINKT}\in\mathsf P$ if there exists some polynomial $\tau$ such that $\mathsf{Gap}_{\tau}\mathsf{MINKT}\in\mathsf P$.
The problem $\mathsf{GapMINKT}$ can be equivalently formulated as follows.
\paragraph{Fact 3.4 ([35]).}
The following are equivalent.
\begin{enumerate}
\item $\mathsf{GapMINKT}\in\mathsf P$.
\item There exist a polynomial-time algorithm $\widetilde K$ and a polynomial $p$ such that
\[K^{p(t)}(x)-\log p(t)\leq\widetilde K(x,1^t)\leq K^t(x)\]
for every string $x\in\{0,1\}^*$ and every integer $t\geq|x|$.
\end{enumerate}
\paragraph{Definition 1.1.}
Symmetry of information for time-bounded Kolmogorov complexity (SoI) refers to the following hypothesis: There exists a polynomial $p$ such that for any strings $x\in\{0,1\}^*$ and $y\in\{0,1\}^*$, for every $t\geq|x|+|y|$,
\[K^{p(t)}(x\mid y)+K^{p(t)}(y)-\log p(t)\leq K^t(x,y).\tag{SoI}\]
Throughout this paper, we refer to this hypothesis as SoI.
The statement $\mathsf{DistNP}\subseteq\mathsf{AvgP}$ means that for every language $L$ in $\mathsf{NP}$ and for every polynomial-time samplable distribution $D$, there exists an errorless heuristic scheme that solves $L$ on inputs sampled from $D$; we refer the reader to the survey of Bogdanov and Trevisan [16] for background on average-case complexity theory.

\paragraph{Prior tools from Section 2.1.}
Our proof of Theorem 1.2 is based on the meta-complexity of time-bounded Kolmogorov complexity: We crucially use the fact that time-bounded Kolmogorov complexity can be approximated in Heuristica.
\paragraph{Lemma 2.1 ([31,33]).}
If $\mathsf{DistNP}\subseteq\mathsf{AvgP}$, then $\mathsf{GapMINKT}\in\mathsf P$.
The second lemma is the reconstruction property of a $k$-wise direct product generator [34], which turned out to be a fundamental tool for analyzing Kolmogorov complexity [33,65,32,35]. A $k$-wise direct product generator $\mathsf{DP}_k:\{0,1\}^n\times\{0,1\}^{nk}\to\{0,1\}^{nk+k}$ is defined as follows:
\[\mathsf{DP}_k(x;z):=(z^1,\ldots,z^k,\langle z^1,x\rangle,\ldots,\langle z^k,x\rangle)\]
for every $x\in\{0,1\}^n$ and every $z=(z^1,\ldots,z^k)\in(\{0,1\}^n)^k$, where $\langle x,y\rangle$ denotes the inner product of $x$ and $y\in\{0,1\}^n$ over $\mathrm{GF}(2)$, i.e., $\langle x,y\rangle:=(\sum_{i=1}^nx_iy_i)\bmod2$. The $k$-wise direct product generator $\mathsf{DP}_k(x;-)$ is a pseudorandom generator secure against an algorithm $D$ if $K(x\mid D)>k+O(\log|x|)$. More formally, we have the following property:
\paragraph{Lemma 2.2 (Deterministic Reconstruction for $\mathsf{DP}_k$; see [35]).}
Assume that $\mathsf E\not\subseteq\mathsf{i.o.SIZE}(2^{\epsilon n})$ for some constant $\epsilon>0$. Then, there exists a polynomial $p$ such that, for every $n\in\mathbb N$, $x\in\{0,1\}^n$, parameters $k,\epsilon^{-1},s\in\mathbb N$, and for every randomized circuit $D$ of size $s$ such that
\[\left|\Pr_{z,r}[D(\mathsf{DP}_k(x;z);r)=1]-\Pr_{w,r}[D(w;r)=1]\right|\geq\epsilon,\]
where $z\sim\{0,1\}^{nk}$, $w\sim\{0,1\}^{nk+k}$, and $r\sim\{0,1\}^s$, it holds that
\[K^{p(ns/\epsilon)}(x\mid D)\leq k+\log p(ns/\epsilon).\]
The assumption of Lemma 2.2 is satisfied in Heuristica, as Buhrman et al. [18] showed that $\mathsf E\not\subseteq\mathsf{i.o.SIZE}(2^{\epsilon n})$ for some constant $\epsilon>0$ if $\mathsf{DistNP}\subseteq\mathsf{AvgP}$. We mention in passing that the proof of Lemma 2.1 also uses Lemma 2.2.

\paragraph{Theorem 4.1.}
If $\mathsf{GapMINKT}\in\mathsf P$ and $\mathsf E\not\subseteq\mathsf{i.o.SIZE}(2^{\epsilon n})$ for some constant $\epsilon>0$, then SoI holds.
\paragraph{Theorem 1.2.}
If $\mathsf{DistNP}\subseteq\mathsf{AvgP}$, then SoI holds.

\paragraph{4. Symmetry of Information from Average-Case Easiness of NP.}
In this section, we present a formal proof of Theorem 1.2. In fact, under a plausible assumption that $\mathsf E$ requires exponential-size circuits to compute, we prove that SoI follows from an approximation algorithm for time-bounded Kolmogorov complexity.
Theorem 4.1 immediately implies Theorem 1.2.
\begin{proof}[Proof of Theorem 1.2]
It is known that the two hypotheses of Theorem 4.1 are implied by the assumption that $\mathsf{NP}$ is easy on average: Buhrman et al. [18] showed that $\mathsf E\not\subseteq\bigcap_{\epsilon>0}\mathsf{i.o.SIZE}(2^{\epsilon n})$ if $\mathsf{DistNP}\subseteq\mathsf{AvgP}$; We also have $\mathsf{GapMINKT}\in\mathsf P$ if $\mathsf{DistNP}\subseteq\mathsf{AvgP}$ by Lemma 2.1.
\end{proof}
We now present the formal proof of Theorem 4.1.
\begin{proof}[Proof of Theorem 4.1]
Let $\widetilde K$ be the polynomial-time algorithm of Fact 3.4 which satisfies the property that there exists a polynomial $p$ such that for every $x\in\{0,1\}^*$ and every $t\geq|x|$,
\[K^{p(t)}(x)-\log p(t)\leq\widetilde K(x;1^t)\leq K^t(x).\tag{9}\]
Fix strings $x\in\{0,1\}^n$ and $y\in\{0,1\}^m$ and an integer $t\geq n+m$. The proof of SoI is given by analyzing the following three values for $z\sim\{0,1\}^{nk}$, $w\sim\{0,1\}^{nk+k}$, $z'\sim\{0,1\}^{m\ell}$, and $w'\sim\{0,1\}^{m\ell+\ell}$:
\[\begin{gathered}
\widetilde K(\mathsf{DP}_k(x;z),\mathsf{DP}_\ell(y;z');1^{t'}),\\
\widetilde K(w,\mathsf{DP}_\ell(y;z');1^{t'}),\\
\widetilde K(w,w';1^{t'}),
\end{gathered}\]
where $t'=t^{O(1)}$, $k\approx K^t(x,y)-\ell$, and $\ell\approx K^{\operatorname{poly}(t)}(y)$ are parameters chosen later.
First, observe that Fact 3.2 implies that $K(w,w')\geq|w|+|w'|-2$ with probability at least $3/4$. Let $\theta:=|w|+|w'|-2-\log p(t')$; using Equation (9), we obtain
\[\Pr_{w,w'}[\widetilde K(w,w';1^{t'})\geq\theta]\geq\frac34.\tag{10}\]
Next, we set the parameter $\ell$ to $K^{p'(t)}(y)-\log p'(t)-1$, where $p'$ is some large polynomial. Consider a randomized circuit $D$ that takes $w'$ as input as well as random bits $w$ and outputs 1 if and only if $\widetilde K(w,w';1^{t'})\geq\theta$. By the contrapositive of Lemma 2.2, $\mathsf{DP}_\ell(y;-)$ is a pseudorandom generator secure against $D$; i.e., $D$ cannot distinguish $\mathsf{DP}_\ell(y;z')$ and $w'$ in the sense that
\[\left|\Pr_{w,z'}[\widetilde K(w,\mathsf{DP}_\ell(y;z');1^{t'})\geq\theta]-\Pr_{w,w'}[\widetilde K(w,w';1^{t'})\geq\theta]\right|<\frac14,\]
which, together with Equation (10), implies that
\[\Pr_{w,z'}[\widetilde K(w,\mathsf{DP}_\ell(y;z');1^{t'})\geq\theta]\geq\frac12.\]
Finally, we compare $\widetilde K(w,\mathsf{DP}_\ell(y;z');1^{t'})$ with $\widetilde K(\mathsf{DP}_k(x;z),\mathsf{DP}_\ell(y;z');1^{t'})$. On one hand, since $|w|=|z|+k$ and $|w'|=|z'|+\ell$, we have
\[\Pr_{w,z'}[\widetilde K(w,\mathsf{DP}_\ell(y;z');1^{t'})\geq|z|+|z'|+k+\ell-2-\log p(t')]\geq\frac12.\tag{11}\]
On the other hand, observe that for some $t':=\operatorname{poly}(t)$,
\begin{align*}
\widetilde K(\mathsf{DP}_k(x;z),\mathsf{DP}_\ell(y;z');1^{t'})
&\leq K^{t'}(\mathsf{DP}_k(x;z),\mathsf{DP}_\ell(y;z'))\\
&\leq K^t(x,y)+|z|+|z'|+O(\log n)
\end{align*}
holds because the strings $\mathsf{DP}_k(x;z)$ and $\mathsf{DP}_\ell(y;z')$ can be computed from $k,\ell,z,z'$, and a program of size $K^t(x,y)$ that outputs $(x,y)$ in time $t$. We now set $k:=K^t(x,y)-\ell+O(\log t)$ so that
\[\Pr_{z,z'}[\widetilde K(\mathsf{DP}_k(x;z),\mathsf{DP}_\ell(y;z');1^{t'})<|z|+|z'|+k+\ell-2-\log p(t')]=1.\tag{12}\]
Let $D_y$ be a randomized circuit that takes an input $w$ and random bits $z'$ and outputs 1 if and only if $\widetilde K(w,\mathsf{DP}_\ell(y;z');1^{t'})<|z|+|z'|+k+\ell-2-\log p(t')$. It follows from Equations (11) and (12) that
\[\Pr_{z,z'}[D_y(\mathsf{DP}_k(x;z);z')=1]-\Pr_{w,z'}[D_y(w;z')=1]\geq1-\frac12=\frac12.\]
Using Lemma 2.2, we obtain
\[K^{\operatorname{poly}(t)}(x\mid D_y)\leq k+O(\log t)=K^t(x,y)-K^{p'(t)}(y)+O(\log t).\]
It follows that for some large polynomial $q$,
\begin{align*}
K^{q(t)}(x\mid y)&\leq K^t(x\mid D_y)+K^t(D_y\mid y)+O(1)\\
&\leq K^t(x,y)-K^{p'(t)}(y)+O(\log t)\\
&\leq K^t(x,y)-K^{q(t)}(y)+\log q(t)
\end{align*}
as desired.
\end{proof}

\begin{thebibliography}{99}
\bibitem[1]{published-1} Dorit Aharonov and Oded Regev. Lattice problems in NP $\cap$ coNP. J. ACM, 52(5):749–765, 2005. doi:10.1145/1089023.1089025.
\bibitem[2]{published-2} Miklós Ajtai. Generating Hard Instances of Lattice Problems (Extended Abstract). In Proceedings of the Symposium on the Theory of Computing (STOC), pages 99–108, 1996. doi:10.1145/237814.237838.
\bibitem[3]{published-3} Eric Allender, Harry Buhrman, Michal Koucký, Dieter van Melkebeek, and Detlef Ronneburger. Power from Random Strings. SIAM J. Comput., 35(6):1467–1493, 2006. doi:10.1137/ 050628994.
\bibitem[4]{published-4} Eric Allender, Mahdi Cheraghchi, Dimitrios Myrisiotis, Harsha Tirumala, and Ilya Volkovich. One-Way Functions and a Conditional Variant of MKTP. In Proceedings of the Foundations of Software Technology and Theoretical Computer Science (FSTTCS), pages 7:1–7:19, 2021. doi:10.4230/LIPIcs.FSTTCS.2021.7.
\bibitem[5]{published-5} Eric Allender, John Gouwar, Shuichi Hirahara, and Caleb Robelle. Cryptographic Hard- ness Under Projections for Time-Bounded Kolmogorov Complexity. In Proceedings of the International Symposium on Algorithms and Computation (ISAAC), pages 54:1–54:17, 2021. doi:10.4230/LIPIcs.ISAAC.2021.54.
\bibitem[6]{published-6} Eric Allender, Joshua A. Grochow, Dieter van Melkebeek, Cristopher Moore, and Andrew Morgan. Minimum Circuit Size, Graph Isomorphism, and Related Problems. SIAM J. Comput., 47(4):1339–1372, 2018. doi:10.1137/17M1157970.
\bibitem[7]{published-7} Eric Allender, Lisa Hellerstein, Paul McCabe, Toniann Pitassi, and Michael E. Saks. Minimizing Disjunctive Normal Form Formulas and AC0 Circuits Given a Truth Table. SIAM J. Comput., 38(1):63–84, 2008. doi:10.1137/060664537.
\bibitem[8]{published-8} Eric Allender, Michal Koucký, Detlef Ronneburger, and Sambuddha Roy. The pervasive reach of resource-bounded Kolmogorov complexity in computational complexity theory. J. Comput. Syst. Sci., 77(1):14–40, 2011. doi:10.1016/j.jcss.2010.06.004.
\bibitem[9]{published-9} Luis Antunes, Lance Fortnow, Dieter van Melkebeek, and N. V. Vinodchandran. Computational depth: Concept and applications. Theor. Comput. Sci., 354(3):391–404, 2006. doi:10.1016/ j.tcs.2005.11.033.
\bibitem[10]{published-10} Luis Antunes, Sophie Laplante, Alexandre Pinto, and Liliana C. M. Salvador. Cryptographic Security of Individual Instances. In Proceedings of Information Theoretic Security (ICITS), pages 195–210, 2007. doi:10.1007/978-3-642-10230-1\_17.
\bibitem[11]{published-11} Luis Filipe Coelho Antunes and Lance Fortnow. Worst-Case Running Times for Average-Case Algorithms. In Proceedings of the Conference on Computational Complexity (CCC), pages 298–303, 2009. doi:10.1109/CCC.2009.12.
\bibitem[12]{published-12} Amos Beimel. Secret-Sharing Schemes: A Survey. In The Third International Workshop on Coding and Cryptology (IWCC), pages 11–46, 2011. doi:10.1007/978-3-642-20901-7\_2.
\bibitem[13]{published-13} Josh Cohen Benaloh and Jerry Leichter. Generalized Secret Sharing and Monotone Functions. In Proceedings of the International Cryptology Conference (CRYPTO), pages 27–35, 1988. doi:10.1007/0-387-34799-2\_3.
\bibitem[14]{published-14} George Robert Blakley. Safeguarding cryptographic keys. In International Workshop on Managing Requirements Knowledge (MARK), pages 313–318, 1979. doi:10.1109/MARK.1979. 8817296.
\bibitem[15]{published-15} Manuel Blum and Sampath Kannan. Designing Programs that Check Their Work. J. ACM, 42(1):269–291, 1995. doi:10.1145/200836.200880.
\bibitem[16]{published-16} Andrej Bogdanov and Luca Trevisan. Average-Case Complexity. Foundations and Trends in Theoretical Computer Science, 2(1), 2006. doi:10.1561/0400000004.
\bibitem[17]{published-17} Andrej Bogdanov and Luca Trevisan. On Worst-Case to Average-Case Reductions for NP Problems. SIAM J. Comput., 36(4):1119–1159, 2006. doi:10.1137/S0097539705446974.
\bibitem[18]{published-18} Harry Buhrman, Lance Fortnow, and Aduri Pavan. Some Results on Derandomization. Theory Comput. Syst., 38(2):211–227, 2005. doi:10.1007/s00224-004-1194-y.
\bibitem[19]{published-19} Harry Buhrman and Elvira Mayordomo. An Excursion to the Kolmogorov Random Strings. J. Comput. Syst. Sci., 54(3):393–399, 1997. doi:10.1006/jcss.1997.1484.
\bibitem[20]{published-20} Lijie Chen, Shuichi Hirahara, and Neekon Vafa. Average-case Hardness of NP and PH from Worst-case Fine-grained Assumptions. In Proceedings of the Innovations in Theoretical Computer Science Conference (ITCS), pages 67:1–67:17, 2022.
\bibitem[21]{published-21} Irit Dinur, Prahladh Harsha, and Guy Kindler. Polynomially Low Error PCPs with polyloglog n Queries via Modular Composition. In Proceedings of the Symposium on Theory of Computing (STOC), pages 267–276, 2015. doi:10.1145/2746539.2746630.
\bibitem[22]{published-22} Irit Dinur and Shmuel Safra. On the hardness of approximating label-cover. Inf. Process. Lett., 89(5):247–254, 2004. doi:10.1016/j.ipl.2003.11.007.
\bibitem[23]{published-23} Joan Feigenbaum and Lance Fortnow. Random-Self-Reducibility of Complete Sets. SIAM J. Comput., 22(5):994–1005, 1993. doi:10.1137/0222061.
\bibitem[24]{published-24} Lance Fortnow and Martin Kummer. On Resource-Bounded Instance Complexity. Theor. Comput. Sci., 161(1\&2):123–140, 1996. doi:10.1016/0304-3975(95)00097-6.
\bibitem[25]{published-25} Lance Fortnow, Troy Lee, and Nikolai K. Vereshchagin. Kolmogorov Complexity with Error. In STACS 2006, 23rd Annual Symposium on Theoretical Aspects of Computer Science, Marseille, France, February 23-25, 2006, Proceedings, pages 137–148, 2006. doi:10.1007/11672142\_10.
\bibitem[26]{published-26} Lance Fortnow and Nick Reingold. PP is Closed Under Truth-Table Reductions. Inf. Comput., 124(1):1–6, 1996. doi:10.1006/inco.1996.0001.
\bibitem[27]{published-27} Eran Gat and Shafi Goldwasser. Probabilistic Search Algorithms with Unique Answers and Their Cryptographic Applications. Electron. Colloquium Comput. Complex., page 136, 2011. URL: https://eccc.weizmann.ac.il/report/2011/136.
\bibitem[28]{published-28} Halley Goldberg and Valentine Kabanets. A Simpler Proof of the Worst-Case to Average-Case Reduction for Polynomial Hierarchy via Symmetry of Information. Electronic Colloquium on Computational Complexity (ECCC), 007, 2022.
\bibitem[29]{published-29} Oded Goldreich and Shafi Goldwasser. On the Limits of Nonapproximability of Lattice Problems. J. Comput. Syst. Sci., 60(3):540–563, 2000. doi:10.1006/jcss.1999.1686.
\bibitem[30]{published-30} Ishay Haviv and Oded Regev. Tensor-based Hardness of the Shortest Vector Problem to within Almost Polynomial Factors. Theory of Computing, 8(1):513–531, 2012. doi:10.4086/ toc.2012.v008a023.
\bibitem[31]{published-31} Shuichi Hirahara. Non-black-box Worst-case to Average-case Reductions within NP. In Proceedings of the Symposium on Foundations of Computer Science (FOCS), pages 247–258, 2018.
\bibitem[32]{published-32} Shuichi Hirahara. Characterizing Average-Case Complexity of PH by Worst-Case Meta- Complexity. In Proceedings of the Symposium on Foundations of Computer Science (FOCS), pages 50–60, 2020.
\bibitem[33]{published-33} Shuichi Hirahara. Non-Disjoint Promise Problems from Meta-Computational View of Pseudor- andom Generator Constructions. In Proceedings of the Computational Complexity Conference (CCC), pages 20:1–20:47, 2020. doi:10.4230/LIPIcs.CCC.2020.20.
\bibitem[34]{published-34} Shuichi Hirahara. Unexpected hardness results for Kolmogorov complexity under uniform reductions. In Proceedings of the Symposium on Theory of Computing (STOC), pages 1038– 1051, 2020. doi:10.1145/3357713.3384251.
\bibitem[35]{published-35} Shuichi Hirahara. Average-case hardness of NP from exponential worst-case hardness assump- tions. In Proceedings of the Symposium on Theory of Computing (STOC), pages 292–302, 2021. doi:10.1145/3406325.3451065.
\bibitem[36]{published-36} Shuichi Hirahara and Mikito Nanashima. On Worst-Case Learning in Relativized Heuristica. In Proceedings of the Symposium on Foundations of Computer Science (FOCS), 2021.
\bibitem[37]{published-37} Shuichi Hirahara, Igor Carboni Oliveira, and Rahul Santhanam. NP-hardness of Minimum Circuit Size Problem for OR-AND-MOD Circuits. In Proceedings of the Computational Complexity Conference (CCC), pages 5:1–5:31, 2018. doi:10.4230/LIPIcs.CCC.2018.5.
\bibitem[38]{published-38} Shuichi Hirahara and Rahul Santhanam. Errorless versus Error-prone Average-case Complexity. In Proceedings of the Innovations in Theoretical Computer Science Conference (ITCS), pages 38:1–38:23, 2022.
\bibitem[39]{published-39} Shuichi Hirahara and Osamu Watanabe. Limits of Minimum Circuit Size Problem as Oracle. In Proceedings of the Conference on Computational Complexity (CCC), pages 18:1–18:20, 2016. doi:10.4230/LIPIcs.CCC.2016.18.
\bibitem[40]{published-40} Rahul Ilango. Approaching MCSP from Above and Below: Hardness for a Conditional Variant and AC0 [p]. In Proceedings of the Innovations in Theoretical Computer Science Conference (ITCS), pages 34:1–34:26, 2020. doi:10.4230/LIPIcs.ITCS.2020.34.
\bibitem[41]{published-41} Rahul Ilango. Constant Depth Formula and Partial Function Versions of MCSP are Hard. In Proceedings of the Symposium on Foundations of Computer Science (FOCS), pages 424–433, 2020.
\bibitem[42]{published-42} Rahul Ilango, Bruno Loff, and Igor Carboni Oliveira. NP-Hardness of Circuit Minimization for Multi-Output Functions. In Proceedings of the Computational Complexity Conference (CCC), pages 22:1–22:36, 2020. doi:10.4230/LIPIcs.CCC.2020.22.
\bibitem[43]{published-43} Russell Impagliazzo. A Personal View of Average-Case Complexity. In Proceedings of the Structure in Complexity Theory Conference, pages 134–147, 1995. doi:10.1109/SCT.1995. 514853.
\bibitem[44]{published-44} Russell Impagliazzo. Relativized Separations of Worst-Case and Average-Case Complexities for NP. In Proceedings of the Conference on Computational Complexity (CCC), pages 104–114, 2011. doi:10.1109/CCC.2011.34.
\bibitem[45]{published-45} Russell Impagliazzo and Leonid A. Levin. No Better Ways to Generate Hard NP Instances than Picking Uniformly at Random. In Proceedings of the Symposium on Foundations of Computer Science (FOCS), pages 812–821, 1990. doi:10.1109/FSCS.1990.89604.
\bibitem[46]{published-46} Russell Impagliazzo and Avi Wigderson. P = BPP if E Requires Exponential Circuits: Derandomizing the XOR Lemma. In Proceedings of the Symposium on the Theory of Computing (STOC), pages 220–229, 1997. doi:10.1145/258533.258590.
\bibitem[47]{published-47} Valentine Kabanets and Jin-yi Cai. Circuit minimization problem. In Proceedings of the Symposium on Theory of Computing (STOC), pages 73–79, 2000. doi:10.1145/335305. 335314.
\bibitem[48]{published-48} Tarik Kaced. Almost-perfect secret sharing. In Proceedings of the International Symposium on Information Theory (ISIT), pages 1603–1607, 2011. doi:10.1109/ISIT.2011.6033816.
\bibitem[49]{published-49} Mauricio Karchmer and Avi Wigderson. On Span Programs. In Proceedings of the Structure in Complexity Theory Conference, pages 102–111, 1993. doi:10.1109/SCT.1993.336536.
\bibitem[50]{published-50} Ker-I Ko. On the Complexity of Learning Minimum Time-Bounded Turing Machines. SIAM J. Comput., 20(5):962–986, 1991. doi:10.1137/0220059.
\bibitem[51]{published-51} Andrei N Kolmogorov. Three approaches to the quantitative definition of information. Problems of information transmission, 1(1):1–7, 1965.
\bibitem[52]{published-52} Troy Lee and Andrei E. Romashchenko. Resource bounded symmetry of information revisited. Theor. Comput. Sci., 345(2-3):386–405, 2005. doi:10.1016/j.tcs.2005.07.017.
\bibitem[53]{published-53} Leonid A. Levin. The Tale of One-Way Functions. Probl. Inf. Transm., 39(1):92–103, 2003. doi:10.1023/A:1023634616182.
\bibitem[54]{published-54} Leonid A. Levin. How do we succeed in tasks like proving Fermat’s Theorem or predicting the Higgs boson?, 2021. An invited talk given at STOC 2021. The transcript is available at https://www.cs.bu.edu/fac/lnd/expo/stoc21/txt.pdf.
\bibitem[55]{published-55} Ming Li and Paul M. B. Vitányi. An Introduction to Kolmogorov Complexity and Its Applications, Third Edition. Texts in Computer Science. Springer, 2008. doi:10.1007/ 978-0-387-49820-1.
\bibitem[56]{published-56} Yanyi Liu and Rafael Pass. On One-way Functions from NP-Complete Problems. Electron. Colloquium Comput. Complex., 28:59, 2021. URL: https://eccc.weizmann.ac.il/report/ 2021/059.
\bibitem[57]{published-57} Luc Longpré and Sarah Mocas. Symmetry of Information and One-Way Functions. Inf. Process. Lett., 46(2):95–100, 1993. doi:10.1016/0020-0190(93)90204-M.
\bibitem[58]{published-58} Luc Longpré and Osamu Watanabe. On Symmetry of Information and Polynomial Time Invertibility. Inf. Comput., 121(1):14–22, 1995. doi:10.1006/inco.1995.1120.
\bibitem[59]{published-59} Zhenjian Lu and Igor Carboni Oliveira. An Efficient Coding Theorem via Probabilistic Representations and Its Applications. In Proceedings of the International Colloquium on Automata, Languages, and Programming (ICALP), pages 94:1–94:20, 2021. doi:10.4230/ LIPIcs.ICALP.2021.94.
\bibitem[60]{published-60} Daniele Micciancio and Oded Regev. Worst-Case to Average-Case Reductions Based on Gaus- sian Measures. SIAM J. Comput., 37(1):267–302, 2007. doi:10.1137/S0097539705447360.
\bibitem[61]{published-61} Noam Nisan and Avi Wigderson. Hardness vs Randomness. J. Comput. Syst. Sci., 49(2):149– 167, 1994. doi:10.1016/S0022-0000(05)80043-1.
\bibitem[62]{published-62} Igor Carboni Oliveira. Randomness and Intractability in Kolmogorov Complexity. In Proceed- ings of the International Colloquium on Automata, Languages, and Programming (ICALP), pages 32:1–32:14, 2019. doi:10.4230/LIPIcs.ICALP.2019.32.
\bibitem[63]{published-63} Igor Carboni Oliveira and Rahul Santhanam. Pseudodeterministic constructions in subex- ponential time. In Proceedings of the Symposium on Theory of Computing (STOC), pages 665–677, 2017. doi:10.1145/3055399.3055500.
\bibitem[64]{published-64} Ran Raz, Omer Reingold, and Salil P. Vadhan. Extracting all the Randomness and Reducing the Error in Trevisan’s Extractors. J. Comput. Syst. Sci., 65(1):97–128, 2002. doi:10.1006/ jcss.2002.1824.
\bibitem[65]{published-65} Hanlin Ren and Rahul Santhanam. Hardness of KT Characterizes Parallel Cryptography. In Proceedings of the Computational Complexity Conference (CCC), pages 35:1–35:58, 2021. doi:10.4230/LIPIcs.CCC.2021.35.
\bibitem[66]{published-66} Detlef Ronneburger. Kolmogorov Complexity and Derandomization. PhD thesis, Rutgers UniversityDept. of Computer Science Laboratory for Computer Sci. Research Hill Center, NJ, USA, 2004.
\bibitem[67]{published-67} Adi Shamir. How to Share a Secret. Commun. ACM, 22(11):612–613, 1979. doi:10.1145/ 359168.359176.
\bibitem[68]{published-68} Amnon Ta-Shma, Christopher Umans, and David Zuckerman. Lossless Condensers, Un- balanced Expanders, And Extractors. Combinatorica, 27(2):213–240, 2007. doi:10.1007/ s00493-007-0053-2.
\bibitem[69]{published-69} Boris A. Trakhtenbrot. A Survey of Russian Approaches to Perebor (Brute-Force Searches) Algorithms. IEEE Annals of the History of Computing, 6(4):384–400, 1984. doi:10.1109/ MAHC.1984.10036.
\bibitem[70]{published-70} Luca Trevisan. Extractors and pseudorandom generators. J. ACM, 48(4):860–879, 2001. doi:10.1145/502090.502099.
\bibitem[71]{published-71} Luca Trevisan and Salil P. Vadhan. Pseudorandomness and Average-Case Complexity Via Uniform Reductions. Computational Complexity, 16(4):331–364, 2007. doi:10.1007/ s00037-007-0233-x.
\bibitem[72]{published-72} Emanuele Viola. On Constructing Parallel Pseudorandom Generators from One-Way Functions. In Proceedings of the Conference on Computational Complexity (CCC), pages 183–197, 2005. doi:10.1109/CCC.2005.16.
\bibitem[73]{published-73} Emanuele Viola. The complexity of constructing pseudorandom generators from hard functions. Computational Complexity, 13(3-4):147–188, 2005. doi:10.1007/s00037-004-0187-1.
\bibitem[74]{published-74} AK Zvonkin and LA Levin. The complexity of finite objects and the algorithmic concepts of randomness and information. UMN (Russian Math. Surveys), 25(6):83–124, 1970.
\end{thebibliography}
\end{document}
